\documentclass{article}
\usepackage{amsmath,amssymb}
\begin{document}

\section*{Whole-problem audit: Higham Chapter 4 Eq. (4.8), Kahan summation}

Please answer the whole mathematical problem, not only the most recent local
lemma.  Start with a conclusion:

\[
\textbf{Conclusion: }\quad
\begin{cases}
\text{Higham's printed returned-sum claim is correct under hypothesis ...},\\
\text{or Higham's printed returned-sum claim is missing hypothesis ...},\\
\text{or the statement is false under the model described below.}
\end{cases}
\]

Then give explicit inequalities/constants, proof steps, and any counterexample.
Do not respond with a plan.

\subsection*{Source claim}

Higham, \emph{Accuracy and Stability of Numerical Algorithms}, 2nd ed.,
Chapter 4, Section 4.3, Eq. (4.8), states for Kahan compensated summation
roughly:
\[
  \widehat s_n=\sum_{i=1}^n x_i(1+\mu_i),\qquad
  |\mu_i|\le 2u+O(nu^2).
\]
This is a source-facing theorem about the ordinary returned stored sum
\(\widehat s_n=s_n\), not the compensated total \(s_n+e_n\).  The Lean
formalization must choose explicit constants and hypotheses replacing the
big-\(O\), and it must not assume exact correction subtraction unless that is
listed as an extra conditional theorem.

\subsection*{Algorithm/model}

The Kahan step is modeled as
\[
\begin{aligned}
  y &= \mathrm{fl}(x-e),\\
  t &= \mathrm{fl}(s+y),\\
  e'&= \mathrm{fl}((t-s)-y),\\
  s'&=t.
\end{aligned}
\]
The base abstract floating-point model supplies only relative-error witnesses
for each rounded operation:
\[
  \mathrm{fl}(a\circ b)=(a\circ b)(1+\delta),\qquad |\delta|\le u,
\]
when the operation is modeled.  There is also a concrete finite round-to-even
format layer, but exact correction subtraction is not automatic for arbitrary
inputs.

\subsection*{Coefficient recurrence}

For one coefficient step, stored coordinates propagate by
\[
  s' = A s + B e,\qquad e' = C s + D e.
\]
Paired coordinates are
\[
  z=s+e,\qquad q=e,
\]
and
\[
  z'=Tz+Rq,\qquad q'=Cz+Hq,
\]
where
\[
  T=A+C,\qquad B_T=B+D,\qquad R=B_T-T,\qquad H=D-C.
\]
The ordinary returned coefficient is
\[
  r=z-q=s,
\]
so
\[
  r'=z'-q'=A z+(B-A)q=A r+Bq.
\]
For a source coefficient injected at step \(i\),
\[
  z_0=B_i+D_i,\qquad q_0=D_i,\qquad r_0=B_i.
\]

\subsection*{Local facts already formalized}

For every actual Kahan coefficient step, with \(0\le u\le 1/64\), there exist
roundoff deltas \(\delta_{\rm sub},\delta_y\) satisfying
\[
  |\delta_{\rm sub}|\le u,\qquad |\delta_y|\le u,
\]
and
\[
\begin{aligned}
  |T-1| &\le 3u^2,\\
  |B_T-1| &\le 2u+9u^2,\\
  |C| &\le u(1+u)^2,\\
  |H| &\le u(1+u)^2(3+u),\\
  |H+\delta_{\rm sub}| &\le 7u^2,\\
  |R-H-\delta_y| &\le u^2,\\
  |R+\delta_{\rm sub}-\delta_y| &\le 8u^2.
\end{aligned}
\]
From these facts, the Lean development proves the actual source-total
coefficient theorem:
for any actual suffix length \(m\) with \(m u^2\le 1/16\),
\[
  |z_m-1|\le 2u+(9+13m)u^2,\qquad |q_m|\le 3u.
\]
This is about \(z=s+e\), not yet about the returned \(r=s=z-q\).

\subsection*{Formalized conditional positive results}

\begin{enumerate}
\item If every correction subtraction \((t-s)\) in the actual prefix trace is
      exact, then the ordinary returned source coefficient theorem is proved:
      \[
        |r_i-1|\le 2u+O(mu^2).
      \]
\item Concrete finite hypotheses implying that exact-subtraction condition
      have been formalized: direct finite correction subtraction, Sterbenz-like
      assumptions, Ferguson-like assumptions, or FastTwoSum-style certificates.
      These are conditional, not arbitrary-input results.
\item The compensated-total returned theorem for \(s_n+e_n\) is proved
      unconditionally under the abstract model with explicit constants.
\item A conditional affine residual-budget theorem is proved: if the final
      retained correction residual is bounded by \(C n u^2\sum_i |x_i|\), then
      the ordinary returned source theorem follows.
\end{enumerate}

\subsection*{Formalized route eliminations / obstructions}

\begin{enumerate}
\item Bare abstract FPModel does not force the first Kahan step from
      \(s=e=0\) to ingest a nonzero finite input exactly.
\item A biased small-\(u\) abstract model satisfies the relative-error contract
      but violates a fixed \(2u+Cu^2\) ordinary returned source cap for
      \([1,0]\) when \(Cu\le 1/2\).
\item Finite operands alone do not imply finite exact subtraction.
\item Tail FastTwoSum order, tail inclusive Sterbenz, tail inclusive Ferguson,
      and direct tail finite correction subtraction are false for arbitrary
      input order.
\item A concrete \(p=5\) binary finite round-to-even trace on four terms
      returns an ordinary stored sum whose source weights cannot all be bounded
      by the pure first-order radius \(2u=1/16\).  This does not by itself
      refute \(2u+O(nu^2)\), but it warns against a pure \(2u\) theorem.
\item The current input-only affine residual majorant cannot supply a fixed
      second-order residual-budget theorem, even for one exact input.
\item The local facts listed above alone cannot imply
      \[
        |r_m-1|\le 2u+K u^2
      \]
      with fixed \(K\): a symbolic two-step local-facts example gives
      \[
        r_m=(1+u)^3=1+3u+3u^2+u^3.
      \]
\end{enumerate}

\subsection*{What I need from you}

Please answer:

\begin{enumerate}
\item Is Higham Eq. (4.8), interpreted as an arbitrary-input ordinary returned
      stored-sum theorem under only the standard relative-error Kahan model,
      actually correct?
\item If yes, what exact hidden invariant beyond the local facts above restores
      \[
        |r_i-1|\le 2u+O(nu^2)?
      \]
      State the invariant in suffix-friendly form and prove it with explicit
      constants.
\item If the theorem needs extra assumptions, identify the minimal missing
      assumptions.  Are they exact correction subtraction, no-underflow/normal
      range, Sterbenz/FastTwoSum order, a guard digit, \(\sum x_i\neq0\), input
      ordering, or something else?
\item If the statement is false under the model, give a minimal mathematical
      counterexample or explain which already-listed obstruction is enough.
\item If the honest theorem should be corrected, what is the best explicit
      source-facing replacement: ordinary returned leading \(3u\)? leading
      \(5u\)? a compensated-total theorem? a conditional exact-subtraction
      theorem? Give explicit constants and proof route.
\end{enumerate}

Please distinguish carefully between:
\[
\text{ordinary returned }s_n,\qquad
\text{compensated total }s_n+e_n,\qquad
\text{conditional exact-subtraction theorem}.
\]

\end{document}
